\documentclass[a4paper]{article}
\usepackage[pdftex]{graphicx}
\usepackage[utf8]{inputenc}
\usepackage{enumerate}
\usepackage{amssymb}
\usepackage{icomma}
\usepackage{siunitx}
\sisetup{locale=DE} 
\usepackage{tikz}
\usepackage{href-ul}
\hypersetup{
	colorlinks=true,
	linkcolor=blue,
	urlcolor=blue}
\usepackage{geometry}
\geometry{a4paper, top=15mm, left=15mm, right=15mm, bottom=15mm,
	headsep=10mm, footskip=12mm}

\begin{document}
\begin{tikzpicture}(10,3)
	\draw[ultra thick](2,0) --(10,0) -- (10,1.3) --(2,1.3) -- (2,0);
	\draw[fill=black](2,0)-- (0,.6) -- (2,.6) -- (2,0);
	\node at (6,.6) {\large Lösungsblatt von \href{https://www.okuyakl.de}{www.okuyakl.de}};
\end{tikzpicture}
\vspace{0.5 cm}

\noindent
\begin{minipage}{0.5\textwidth}
	\noindent{\bf Aufgabe 1. a)}\\
	\includegraphics[width=7 cm]{tangent052}
\end{minipage}
\hfill
\begin{minipage}{0.5\textwidth}
	\noindent{\bf Aufgabe 1. b)}\\
	Die Ortslinie ist der Thaleskreis um $M$ durch $P$ und $A$.
\end{minipage}

\noindent{\bf Aufgabe 2.}\\
Hierbei handelt es sich um zwei Parallelen, welche ein Paar bilden und von der Mittelparallelen $g$ den Abstand $\SI{8}{\centi\meter}$ haben:
$$M=\{ P| d(P;g)=\SI{8}{\centi\meter}\}$$

\noindent {\bf Aufgabe 3. a)}\\
$$A=\{ P|d(P;M_1)>r_1 \land P| d(P;M_2)<r_2 \} $$

\noindent {\bf Aufgabe 3. b)}\\
$$B=\{ P|d(P;g)\ge d(P;h)\} $$

\noindent {\bf Aufgabe 3. c)}\\
Der Winkel zwischen $w_1$ und $w_2$ ist $\alpha + \beta$.
Für die Winkel $\alpha$ und $\beta$ gilt:
$$
\renewcommand{\arraystretch}{1.5}
\begin{array}{rcll}
2\cdot \alpha + 2 \cdot \beta &=& 180^\circ & |:2 \\ 
               \alpha + \beta &=& 90^\circ  

\end{array}
$$
$w_1$ und $w_2$ schließen also einen rechten Winkel ein.

\noindent {\bf Aufgabe 4. a)}\\
$M_1$ sind zwei Ortslinien, und zwar die  Winkelhalbierenden von $g_1$ und $g_2$:
$$M_1=\{P|d(P;g_1)=d(P;g_2)\}$$
 
\noindent {\bf Aufgabe 4. b)}\\
$M_2$ ist eine Ortslinie, die Mittelparallele zwischen $h_1$ und $h_2$:
$$M_2=\{P|d(P;h_1)=d(P;h_2)\}$$

\noindent {\bf Aufgabe 4. b)}\\
$K$ ist ein Ortsbereich, nämlich das Innere der Kreises $k(M;r)$ um $M$ mit dem Radius $r$:  
$$K=\{P|d(P;M)\le r\}$$

\noindent {\bf Aufgabe 5.)}\\

\begin{minipage}{0.5\textwidth}
\includegraphics[width=5 cm]{papa160}
\end{minipage}
\hfill
\begin{minipage}{0.5\textwidth}
P liegt auf einem Paar paralleler Geraden im Abstand von jeweils $\SI{1,5}{\centi\meter}$ von $g$
\end{minipage}

\noindent {\bf Aufgabe 6.a)}\\
Die Punktmenge ist eine Ortslinie, weil der Abstand immer einen bestimmten Wert hat. 
Sie beschreibt eine Kreislinie um $M$ mit dem Radius $\SI{2,8}{\centi\meter}$.
  
\noindent {\bf Aufgabe 6.b)}\\
Diese Punktmenge ist ein Ortsbereich, denn der Abstand ist größer oder gleich $r_2$. Sie 
beschreibt den Bereich außerhalb eines Kreises um $M$ mit dem Radius $r_2$.

\noindent {\bf Aufgabe 7.)}\\
\begin{minipage}{0.6\textwidth}
\includegraphics[width=6 cm]{pum160}
\end{minipage}
\hfill
\begin{minipage}{0.4\textwidth}
{\bf a)} P liegt innerhalb eines Thaleskreises zwischen $A$ und $C$\\
{\bf b)} P' liegt auf und rechts neben der Mittelsenkrechten von $A$ und $B$. 
\end{minipage}

\noindent {\bf Aufgabe 8.)}\\
$$G=\{ P| \SI{2}{\centi\meter} \le d(P;M) < \SI{4}{\centi\meter} \}$$

\begin{center}
	\includegraphics[width=7 cm]{../../viecher/endcomic.pdf}
	
	Hier geht es zurück zum \href{https://www.okuyakl.de/math/m8ortL052/aa052.pdf}{Aufgabenblatt}
\end{center}

\end{document}

